\chapter{Estado del Arte}
\label{chap:estado-arte}

Este capítulo desarrolla la base conceptual necesaria para abordar las preguntas de investigación de la sección~\ref{sec:preguntas-investigacion}. La primera parte examina la discusión como proceso de aprendizaje, las funciones que desempeña el instructor al facilitarla y la información disponible cuando el intercambio ocurre por escrito y en momentos distintos. La segunda estudia cómo los sistemas inteligentes en educación han representado algunas de esas funciones, desde los agentes pedagógicos hasta los sistemas basados en modelos de lenguaje de gran escala y las arquitecturas multiagente.

\section{Facilitación de discusiones académicas}

\subsection{La discusión como proceso de aprendizaje}

La primera pregunta de la fase exploratoria plantea qué diferencia una discusión académica de otras formas de aprendizaje. La literatura la describe como un proceso de construcción conjunta en el que los participantes formulan ideas, las contrastan y las revisan a partir de las aportaciones del grupo \parencite{Murphy2018,Mercer2000,Glina2012}. Su valor aparece cuando una contribución retoma, cuestiona o amplía otra, ya que esta relación entre aportaciones permite observar el razonamiento del grupo y ofrece oportunidades de participación distintas de las de una clase expositiva.

Para diseñar el sistema era necesario pasar de esa descripción general a decisiones observables: qué señales lee un instructor, cuándo espera y qué acción realiza si decide intervenir.

\subsection{El rol del instructor en la facilitación}

La segunda pregunta examina qué hace el instructor cuando facilita una discusión. La literatura muestra que lee el hilo de discusión, interpreta si los participantes exploran, integran ideas o llegan a conclusiones, y decide si conviene intervenir o esperar \parencite{Abdous2011,Rovai2007,GarrisonEtAl2000}. Cuando interviene, la literatura distingue de forma recurrente funciones organizacionales, intelectuales y sociales \parencite{Paulsen1995,Abdous2011,BlignautTrollip2003}. Estas funciones proporcionan la base de los roles utilizados por el sistema:

\begin{itemize}
    \item \textbf{Rol organizacional}: gestiona la estructura de la discusión, las normas de participación y la continuidad del hilo de discusión \parencite{Paulsen1995,Abdous2011,PilkingtonWalker2003}.
    \item \textbf{Rol intelectual}: formula preguntas, conecta argumentos y orienta el razonamiento del grupo \parencite{Abdous2011,Murphy2018,PilkingtonWalker2003}.
    \item \textbf{Rol social}: favorece la participación y la relación entre las contribuciones de los estudiantes \parencite{Paulsen1995,Abdous2011,BlignautTrollip2003}.
    \item \textbf{Rol afectivo}: reúne acciones de motivación, reconocimiento y apoyo emocional documentadas en los agentes pedagógicos \parencite{SikstromEtAl2022}.
    \item \textbf{Rol moderador}: identifica contenido problemático y contempla su escalado para revisión \parencite{Rovai2007,KorreEtAl2025}.
\end{itemize}

Los tres primeros roles trasladan categorías recurrentes de la literatura. Durante el diseño, las funciones afectivas y de moderación se separaron en dos roles operativos porque requieren instrucciones y herramientas diferentes. Esta separación constituye una decisión de implementación cuya utilidad pedagógica todavía no se ha validado.

\subsection{Discusiones asíncronas y comunicación mediada}

La comunicación mediada por computador comprende los intercambios que se desarrollan mediante sistemas digitales. Dentro de esta categoría, este trabajo se centra en discusiones escritas y asíncronas, en las que los participantes intervienen en momentos distintos \parencite{Rovai2007,MasseyEtAl2019}. La tercera pregunta de la fase exploratoria examina qué información puede utilizar el facilitador en estas condiciones. La asincronía ofrece tiempo para elaborar una respuesta y permite participar sin coincidir en horario, pero elimina señales inmediatas de comprensión, desacuerdo o retirada. El instructor tiene que reconstruir la trayectoria a partir del texto, las relaciones entre mensajes y el tiempo transcurrido. La ausencia de una respuesta puede indicar reflexión, desconexión o cierre; el volumen de mensajes tampoco garantiza avance conceptual \parencite{Rovai2007,MasseyEtAl2019}.

Esta ambigüedad se combina con un problema de atención. Cada hilo conserva su historial, pero el instructor debe seguir muchos intercambios en paralelo y decidir cuáles requieren presencia. La dificultad aumenta en cursos masivos, donde el volumen de participantes hace más costoso sostener una facilitación continua \parencite{GuEtAl2021,YuEtAl2024maic}. Por eso, el problema técnico comprende la interpretación de la trayectoria del hilo de discusión, la decisión de actuar y, cuando corresponde, la generación de una respuesta.

\section{Sistemas inteligentes en educación}

Esta sección examina los sistemas inteligentes en educación para identificar conceptos y patrones de diseño relacionados con el sistema propuesto. Primero se presentan los agentes pedagógicos y los sistemas de tutoría inteligente. Después se estudian los agentes basados en modelos de lenguaje de gran escala, los sistemas multiagente y las aplicaciones desplegadas a mayor escala. Cuando los trabajos revisados incluyen una evaluación, se describe únicamente la evidencia necesaria para interpretar su relación con este trabajo.

\subsection{Agentes pedagógicos y comunicativos}

Los agentes pedagógicos son sistemas de software que se comunican con estudiantes en entornos de aprendizaje digital. El marco de referencia más empleado para evaluarlos es la Comunidad de Indagación \parencite{GarrisonEtAl2000}, que organiza el aprendizaje en línea en torno a tres presencias. La presencia cognitiva abarca la construcción de conocimiento; la social, la proyección de identidad y la cohesión grupal; y la docente, el diseño y la facilitación. Las funciones comunicativas documentadas en estos agentes se dividen en procesos intrapersonales (motivación, autorregulación), andamiaje interpersonal (retroalimentación, pistas) y facilitación de la colaboración grupal \parencite{SikstromEtAl2022}. Los sistemas actuales son efectivos en los dos primeros tipos de función, pero aún carecen de capacidad para sostener conversaciones recíprocas fluidas a nivel de grupo \parencite{SikstromEtAl2022}. La evidencia reciente muestra que las arquitecturas con múltiples agentes y roles diferenciados (tutor, par, agente enseñable) producen beneficios pedagógicos que los sistemas de un solo agente no logran \parencite{LippertEtAl2020}, lo que sugiere que la diferenciación de roles es una vía para resolver esa brecha.

En el contexto específico de las discusiones mediadas por computador, la investigación ha documentado qué roles de facilitación sostienen conversaciones productivas. Los roles de gestión permiten mantener el foco, gestionar tiempos y subtareas; los de construcción de comunidad establecen normas de participación y crean un espacio donde compartir ideas sin riesgo; y los de argumentación consisten en formular, cuestionar y contracuestionar razonamientos con justificación explícita \parencite{PilkingtonWalker2003}. La síntesis de \textcite{Abdous2011} organiza las funciones del facilitador en categorías intelectual, social y de gestión, e incluye el resumen de las contribuciones entre las habilidades relevantes para sostener la discusión. Estas acciones de gestión, construcción de comunidad y argumentación proporcionan antecedentes concretos para definir las intervenciones que genera el sistema según el estado de la discusión.

La forma más estudiada de agente pedagógico adaptativo es el sistema de tutoría inteligente (ITS). Una revisión de 85 estudios muestra que los ITS producen un tamaño del efecto de $d \approx 0{,}76$ (d de Cohen, que mide la diferencia entre grupos en unidades de desviación típica; 0,76 indica una mejora sustancial) respecto a la instrucción en clase, comparable al de la tutoría humana, con sistemas como AutoTutor y Cognitive Tutor de Carnegie Learning como referentes empíricos \parencite{VanLehn2011}. Los sistemas más avanzados alcanzan fiabilidad comparable a la de evaluadores humanos en el rastreo del dominio de contenidos y las características cognitivas de los estudiantes \parencite{GraesserEtAl2017}. Sin embargo, los ITS asumen un modelo de dominio con respuestas correctas e incorrectas, supuesto que no se sostiene en discusiones académicas abiertas donde puede no existir una única respuesta de referencia. La evidencia sobre aprendizaje productivo por fracaso añade una restricción relevante, pues los estudiantes que trabajan un problema sin asistencia antes de recibir instrucción obtienen mejores resultados de transferencia que los que reciben instrucción desde el inicio \parencite{Kapur2016}, lo que indica que intervenir demasiado pronto puede interferir con el aprendizaje.

De esta revisión se extraen dos aportaciones directas para el diseño del sistema. La taxonomía de roles de \textcite{PilkingtonWalker2003} (gestión, construcción de comunidad y argumentación) define los tipos de intervención que el sistema debe producir, y el principio de aprendizaje productivo por fracaso \parencite{Kapur2016} sustenta la regla de no intervenir mientras la dinámica de la discusión no ha alcanzado un punto de bloqueo.

\subsection{Modelos de lenguaje de gran escala y agentes en educación}

Los modelos de lenguaje de gran escala pueden realizar una tarea a partir de instrucciones y ejemplos incluidos en la entrada. \textcite{BrownEtAl2020} evaluaron esta capacidad sin proporcionar ejemplos, con un solo ejemplo y con varios ejemplos, sin actualizar los parámetros del modelo para cada tarea. Esta propiedad permite utilizar un mismo modelo para distintas operaciones lingüísticas mediante cambios en las instrucciones y el contexto que recibe.

El término ingeniería del contexto se ha propuesto para describir el diseño, la estructuración y la gestión del entorno informativo que recibe un agente basado en LLMs \parencite{Vishnyakova2026}. El informe técnico de la Cátedra iDanae (UPM y Management Solutions) lo concreta como la optimización de la cantidad de información introducida en una llamada al modelo y de la forma en que se presenta \parencite{IdanaeUPM2026}.

La capacidad de seguir instrucciones y ejemplos permite utilizar modelos generales en distintas tareas lingüísticas. Al mismo tiempo, sus respuestas pueden contener información no sustentada, sesgos o problemas de privacidad \parencite{ChuEtAl2025}. Por esta razón, el sistema conserva sus salidas como respuestas candidatas que pueden inspeccionarse, y su validez pedagógica queda pendiente de una evaluación posterior.

La aplicación de estos agentes en educación se ha concentrado en tutoría personalizada y diseño curricular, con escasa evidencia sobre facilitación de discusiones grupales. Las revisiones sistemáticas disponibles muestran que estas aplicaciones emplean memoria, planificación y herramientas \parencite{ChuEtAl2025}, pero la mayoría se sitúa en niveles intermedios de madurez tecnológica (TRL, \textit{Technology Readiness Level}, 4-6) y la evidencia longitudinal sobre su impacto en el aprendizaje sigue siendo limitada \parencite{KostopoulosEtAl2025}. Alrededor del 70\% de los estudios asignan a la IA la función de estructurar condiciones previas de aprendizaje, mientras que la facilitación activa de la dinámica del grupo aparece con menor frecuencia \parencite{CasebourneEtAl2025}.

Entre los trabajos sobre facilitación activa, \textcite{ChenEtAl2023tutoring} organizan la tutoría adaptativa en tres procesos encadenados de interacción, reflexión y reacción. Otro antecedente es el tutor socrático basado en LLM estudiado por \textcite{DegenAsanov2025}, que aplica la taxonomía de preguntas de \textcite{Paul1990}. En un estudio con 65 estudiantes, este tutor ofreció mayor apoyo al pensamiento crítico que los métodos comparados. En conjunto, ambos trabajos muestran formas de distribuir una intervención educativa entre procesos o responsabilidades diferenciadas, por lo que constituyen antecedentes para la arquitectura que se presenta en el capítulo siguiente.

\subsection{Intervenciones en discusiones académicas}

La investigación sobre intervenciones automatizadas en discusiones complementa ese trabajo orientado a agentes al analizar directamente el problema de cuándo actuar, qué producir y a quién dirigirse.

La investigación en esta área muestra que la generación de intervenciones es técnicamente viable, y que el quién y el cuándo son tan determinantes como el qué. El trabajo existente se ha centrado principalmente en tareas de clasificación, como la detección de toxicidad y el análisis de postura; la producción de intervenciones destinadas a mejorar la calidad de la discusión recibe menos atención \parencite{KorreEtAl2025}. Los estudios de revisión sobre moderación asistida por IA también muestran que el campo se concentra en técnicas restrictivas, como la eliminación de contenido, con escasa exploración de los impactos a nivel grupal \parencite{YuEtAl2024review}. Sobre el cuándo, la evidencia muestra que el descarrilamiento conversacional puede predecirse antes de que ocurra siguiendo la trayectoria del tono a lo largo del hilo \parencite{ChangDanescu2019}. Sobre el qué, hay evidencia de que la generación automatizada de sugerencias de facilitación es factible, con una precisión de 0,679 y una exhaustividad (\textit{recall}) de 0,699 sobre 1.600 publicaciones reales, aunque la verificación humana sigue siendo necesaria para tareas complejas \parencite{GuEtAl2021}. Sobre el quién, un chatbot moderador en discusiones deliberativas (n=64) muestra que dirigirse específicamente a participantes reticentes produce mayor alineación de opiniones, distribución más equitativa de contribuciones y mayor cohesión grupal que intervenir sobre el hilo en general \parencite{KimEtAl2021}.

Estas tres líneas de evidencia informan directamente el diseño del sistema. La predictibilidad del descarrilamiento \parencite{ChangDanescu2019} fundamenta el módulo de clasificación de estados de la discusión, la viabilidad de la generación automatizada \parencite{GuEtAl2021} respalda como antecedente el enfoque técnico adoptado, y la efectividad de las intervenciones dirigidas a participantes específicos \parencite{KimEtAl2021} orienta el diseño de las intervenciones de tipo social.

\subsection{Sistemas multiagente en el aula: SimClass y MAIC}

SimClass\footnote{\url{https://github.com/THU-MAIC/SimClass}} y MAIC\footnote{El repositorio actual del proyecto puede consultarse en \url{https://github.com/THU-MAIC/OpenMAIC}.} fueron desarrollados por equipos de la Universidad Tsinghua, comparten varios autores y estudian arquitecturas multiagente aplicadas a la educación. Cada uno examina una situación distinta y ambos permiten estudiar cómo se distribuyen funciones educativas entre agentes especializados.

SimClass estudia si agentes LLM con roles diferenciados pueden reproducir dinámicas sociales emergentes en una simulación de aula. La base arquitectónica son los agentes generativos dotados de flujo de memoria, reflexión y planificación, que producen comportamientos sociales emergentes en entornos de simulación \parencite{ParkEtAl2023}. SimClass aplica este enfoque a un aula virtual donde un agente instructor, un agente asistente y varios agentes estudiante interactúan con personas diferenciadas (Class Clown, Deep Thinker, Note Taker, Inquisitive Mind) \parencite{ZhangEtAl2024simclass}. El sistema se evalúa con el Sistema de Análisis de Interacción de Flanders (FIAS) y el marco de la Comunidad de Indagación; en experimentos con más de 400 estudiantes de posgrado, produjo comportamientos de grupo emergentes (enseñanza colaborativa, acompañamiento emocional, gestión de la disciplina) y los estudiantes reportaron alta presencia e interacción efectiva \parencite{ZhangEtAl2024simclass}. A partir de estos resultados, SimClass aporta evidencia sobre el uso de roles diferenciados y la aparición de dinámicas grupales en un aula simulada. Sin embargo, su unidad de interacción es el aula virtual completa, mientras que este trabajo analiza hilos de discusión asíncronos recuperados de una fuente externa.

MAIC estudia cómo distribuir funciones docentes entre agentes LLM en una plataforma masiva. El sistema opera en tres etapas (Read, transformación de materiales del curso; Teach, generación de guiones y acciones de enseñanza; y Learn, interacción en el aula) con cinco roles de agente (Teacher, Companion, Assistant, Analyzer y Manager) \parencite{YuEtAl2024maic}. En experimentos con más de 500 estudiantes y más de 100.000 registros de aprendizaje en la Universidad Tsinghua, los autores identifican correlaciones entre la interacción con el sistema y el rendimiento académico. También señalan que los agentes instructores usaban guiones estandarizados para todos los estudiantes y que el impacto observado sobre pensamiento profundo y calidad de la discusión fue limitado \parencite{YuEtAl2024maic}. Por tanto, MAIC resulta relevante para este trabajo por la distribución de funciones docentes entre agentes y por su aplicación a una plataforma educativa masiva. Esa separación de responsabilidades se traslada aquí al análisis y la facilitación de hilos de discusión asíncronos.

\subsection{Plataformas a escala: MOOCs y aplicaciones empresariales}

Más allá de los prototipos de investigación, la aplicación de IA a la facilitación de discusiones educativas también ha generado soluciones comerciales a escala.

Las aplicaciones existentes utilizan modelos de acceso e integración diferentes. Coursera incorporó asistentes basados en GPT-4 en más de 100 cursos, con tiempos de respuesta un 47\% más rápidos que los tutores humanos \parencite{Coursera2024}, mientras que edX reportó 1,2 millones de estudiantes en seis meses utilizando laboratorios virtuales impulsados por IA \parencite{MarketGrowthReports2026}. CONNECT, desarrollado por Human2Human AI, ofrece sesiones de discusión en grupos pequeños moderadas por un facilitador de IA y se integra con LMS compatibles mediante LTI (\textit{Learning Tools Interoperability}) \parencite{Human2Human2025}. Estas aplicaciones muestran distintas formas de incorporar IA a plataformas educativas. El sistema desarrollado en este trabajo se concentra en intervenciones trazables sobre discusiones asíncronas y en un contrato común para recibir los hilos.

\section{Conclusiones del estado del arte}

Los sistemas inteligentes en educación muestran resultados favorables en tutoría individual y permiten simular aulas con múltiples roles, aunque la facilitación de hilos de discusión asíncronos sigue menos explorada. Según \textcite{CasebourneEtAl2025}, alrededor del 70\% de los estudios revisados asignan a la IA la función de estructurar las condiciones de aprendizaje, mientras que su participación activa en una conversación grupal aparece con menor frecuencia. En este contexto, SimClass estudia dinámicas de aula simuladas y MAIC distribuye funciones docentes en una plataforma educativa masiva. Entre los sistemas revisados, ninguno aborda un problema similar al de este trabajo, que consiste en recibir un hilo de discusión procedente de una fuente externa, analizar si necesita apoyo y producir una intervención cuya generación pueda inspeccionarse.

La facilitación automatizada de un hilo asíncrono requiere dos capacidades. La primera es interpretar el hilo para reconocer señales como estancamiento, desvío o conflicto \parencite{WeiEtAl2017,AgrawalEtAl2015,LiuEtAl2015}. La segunda es producir una respuesta contextualizada, como una pregunta de seguimiento, una conexión entre aportaciones o una síntesis \parencite{SengarEtAl2025,KilincKececioglu2024,BrownEtAl2020}. Una única instrucción de entrada (\textit{prompt}) podría reunir ambas tareas, aunque combinaría sus instrucciones y estados y reduciría la independencia de cada etapa y la inspección de las decisiones intermedias. Por esta razón, el diseño las organiza en etapas sucesivas de clasificación, decisión y generación.

El marco de \textcite{DwivediEtAl2026} proporciona conceptos para describir sistemas que combinan planificación, memoria compartida, herramientas y distribución de tareas entre agentes. Estos conceptos se utilizarán en el capítulo siguiente para caracterizar la arquitectura implementada.

La revisión también produjo un repertorio de técnicas y condiciones para decidir cuándo conviene esperar. El inventario de fuentes está disponible en el archivo de revisión bibliográfica del repositorio\footnote{\href{https://github.com/magrimal/tfm-2026-discussion-moderation/blob/main/docs/literature/papers.csv}{Repositorio del proyecto, \texttt{docs/literature/papers.csv}}}, y cada técnica conserva sus referencias. La evidencia sobre fracaso productivo advierte que una ayuda temprana puede interrumpir una dificultad todavía útil en tareas con una solución identificable \parencite{Kapur2016,VanLehn2011}. Su aplicación a discusiones abiertas constituye una hipótesis de diseño que requiere validación posterior.

El capítulo siguiente presenta la arquitectura del sistema y explica cómo organiza la clasificación, la decisión de intervenir y la generación de la respuesta en etapas sucesivas.
